%% Review response for Mauro De Marchis. Same substance as
%% docs/paper1_internal_review.md, with the internal shorthand resolved:
%% memory keys, script paths and repository-relative links mean nothing to a
%% reader outside the project. Letterhead shared with manuscript/cover_letter.tex.
\documentclass[11pt]{article}
\usepackage[left=2.2cm, right=2.2cm, bottom=1.8cm, top=3.2cm, headheight=3.2cm,
            headsep=0.3cm, footskip=1.5cm]{geometry}
\usepackage[T1]{fontenc}
\usepackage[utf8]{inputenc}
\usepackage{lmodern}
\usepackage{parskip}
\setlength{\parskip}{0.5em}
\usepackage[hidelinks]{hyperref}
\usepackage{enumitem}
\usepackage{titlesec}
\usepackage{graphicx}
\usepackage{booktabs}
\usepackage{tabularx}
\usepackage{fancyhdr}
\usepackage{siunitx}

\titleformat{\section}{\normalfont\bfseries\large}{\thesection.}{0.5em}{}
\titleformat{\subsection}{\normalfont\bfseries}{\thesubsection}{0.5em}{}
\titlespacing*{\section}{0pt}{0.9em}{0.3em}
\titlespacing*{\subsection}{0pt}{0.7em}{0.25em}

\pagestyle{fancy}
\fancyhf{}
\renewcommand{\headrulewidth}{0pt}
\renewcommand{\footrulewidth}{0pt}
\fancyhead[C]{\raggedright\includegraphics[width=0.5\textwidth]{../manuscript/unipa_header.png}}
\fancyfoot[C]{\includegraphics[width=0.9\textwidth]{../manuscript/unipa_footer.png}}

\newcommand{\itemhead}[1]{\vspace{0.3em}\noindent\textbf{#1}\par}

\begin{document}
\thispagestyle{fancy}

\begin{center}
{\large\bfseries Paper 1: response to the internal review}\\[0.3em]
Cicero Martins Jr.\\
\textit{To: Prof.\ Mauro De Marchis} \hfill 13 October 2026
\end{center}

\vspace{0.2em}
\hrule
\vspace{0.5em}

Each of the five observations is recorded below with what was found on
checking it against the manuscript, so that the scope of each change is known
before it is made. None has been implemented yet.

Two further sections follow. The first carries the sections not yet reviewed,
scanned against the criterion the five observations share. The second measures
the manuscript against the journal's stated requirements.

\section{The five observations}

\itemhead{1. Carmelo Nasello and Antonino Maltese as co-authors}

Three details are not settled by the observation and would benefit from a
decision: position in the author order, affiliation where it differs from
Dipartimento di Ingegneria, and the CRediT roles. The three current entries
follow the pattern of the Sentinel-1 reservoir manuscripts, where supervising
co-authors carry Conceptualization, Supervision, Methodology and Writing
(review and editing). Antonino Maltese supplied the classification method that
Section 3.4 builds on, which may warrant a role beyond that pattern.

Elsevier does not permit authorship changes after acceptance, and changes
before acceptance require written agreement from every author, so the list
should be definitive at submission.

\itemhead{2. Lagoon geometry in the Introduction}

The duplication is complete. The second Introduction paragraph and Section 2
carry the same quantities: 11~km north to south, 2.5~km east to west, 2200~ha,
the northern inlet at 400~m wide and 0.30--0.40~m deep, the southern at 2900~m
and 1.0--1.5~m, and shoot densities of 1000--1200 plants~\si{\per\metre\squared}.
The figures can therefore be cut from the Introduction without loss.

Two elements of that paragraph are argument rather than description, and a
decision on where they belong is separate from the geometry: the sub-metre mean
depth, which is the ground for rejecting the depth-averaging assumption, and
the 75\% reduction in \emph{P.\ oceanica} extent over 2003--2024, which
motivates the vegetated arm of the design.

\itemhead{3. Age of the reference set}

Twenty-two of the 29 entries predate 2015, and the Introduction draws on
nineteen works of which most fall in that group.

The set divides between the formulations the model implements and the claims
that frame the problem. Booij et al.\ (1999), Baptist et al.\ (2007), Kernkamp
et al.\ (2011) and Liu and Weisberg (2011) are cited because the model implements
those specific formulations, and a newer reference does not substitute for
them. The candidates for refreshment are the framing claims, where a recent
synthesis can stand beside or in place of an older source, and the
Mediterranean lagoon context presently carried by Niedda and Greppi (2007) and
Umgiesser et al.\ (2014).

A forward-citation search on each older work, filtered to coastal lagoon
hydrodynamics, seagrass canopy flow and Lagrangian validation, is the method
proposed.

\itemhead{4. Figure 1}

The margin is in panel (a), between the map and the axes frame, and the panel
takes its limits from a constant that does not match the mesh it paints:

\begin{center}
\begin{tabular}{llll}
\toprule
 & Panel (a) limits & Mesh extent & Margin \\
\midrule
Longitude & 11.950 to 12.570 & 11.9600 to 12.5581 & 0.0100 W, 0.0119 E \\
Latitude  & 37.690 to 38.072 & 37.7000 to 38.0631 & 0.0100 S, 0.0089 N \\
\bottomrule
\end{tabular}
\end{center}

Panel (b) takes its limits from the lagoon bounding box and paints to the
frame, so no margin appears there. The latitude band in panel (a) is 2.6\% of
the panel range, which accounts for the millimetre-scale difference in apparent
map height. Matching panel (a) to the mesh extent removes it.

\itemhead{5. Justification of the choices in Section 3.1}

Section 3.1 states the mesh, the ten sigma layers, the eight partitions and the
time step without the ground for each. Two can be supplied immediately. The
partition count is a throughput choice rather than a physical one, and the
supporting measurement exists, the eight-rank configuration reaching about nine
hours per nine simulated days on the 64-core machine; saying so removes the
implication that eight carries physical meaning. The discretisation method is
documented in the project record, from the Navier--Stokes equations to the
discrete system the solver advances, and Section 3.1 can draw on it.

The mesh resolution and the layer count are different. Neither has a
convergence test behind it. A defensible treatment is to state the constraint
each choice answers and to record that no sensitivity test was run, rather than
to imply a convergence study that does not exist.

\section{The sections not yet reviewed}

Four of the five observations express one criterion, that a value stated is a
value to be grounded. The remaining sections were scanned against it by rule: a
sentence asserting a setting is flagged where neither it nor its neighbour
carries a causal clause, a constraint answered, or a citation.

Eleven sentences flag in Section 3. Two are grounded on re-reading, the
evaporation sign convention and the sediment thickness taper, which the pattern
missed. Nine assert a setting with no ground:

\begin{center}
\begin{tabular}{lll}
\toprule
Section & Setting & Value \\
\midrule
3.1 & Uniform Manning coefficient & 0.023~s~m$^{-1/3}$ \\
3.1 & External forcing update interval & 300~s \\
3.1 & Courant limit, step cap, growth factor & 0.7, 30~s, 1.1 \\
3.2 & Spectral discretisation & 24 directions, 24 frequencies, 0.05--1.0~Hz \\
3.2 & Source terms & third generation, Komen, quadruplets, refraction \\
3.2 & Processes disabled & triads, wave setup \\
3.2 & Breaking and bottom friction & $\alpha=1.0$, $\gamma=0.73$; JONSWAP 0.067 \\
3.2 & Minimum depth & 0.05~m \\
3.2 & Coupling exchange interval & 3600~s \\
\bottomrule
\end{tabular}
\end{center}

Six of the nine fall in Section 3.2, where the review stopped. The
concentration is not an artefact of the scan: Section 3.2 is a list of SWAN
settings, and the only ground it offers anywhere is for the NetCDF format of
the communication file. Applied to Sections 4 and 5 the same scan returns
reported measurements rather than settings, so the criterion is a Methods one.

\subsection{Points in the same sections that a scan does not reach}

The two nested SWAN grids, at 400~m and 100~m, have no stated relation to the
unstructured mesh they exchange with. Whether the inner grid resolves the
inlets, which are 400~m and 2900~m wide, is the question a reader of Section
3.1 carries into 3.2.

The 3600~s exchange interval sets the time resolution of the whole coupling,
against drifter deployments lasting 0.4 to 7.2 hours.

Stationary mode within each exchange is explained mechanically. The ground for
it being adequate, that the wave field equilibrates faster than the hourly
exchange in a fetch-limited basin, is not given.

Section 3.4 does not say why the vegetation module was used rather than
trachytopes, which is the route the earlier literature on this lagoon takes.
The reason held in the project record is that the trachytope canopy term does
not enter the momentum equation in this build. How much of that belongs in the
manuscript is a judgement, but silence invites the question.

Section 3.6 states the exclusion of the suspended sediment contribution to
density as a property of the design. A reviewer may read it as a parameter
tuned until the runs completed. The evidence against that reading, that the two
affected members abort at the same cell and the same instant to six decimal
places on independent attempts, is not in the manuscript.

\section{The journal's requirements}

Two sources carry assessable criteria: Elsevier's reviewer guidance for
original research articles, and the \textit{Estuarine, Coastal and Shelf
Science} Guide for Authors. The editorial by the journal's Editors-in-Chief,
Borja and Elliott (2024), was read and concerns the shortage of referees rather
than the assessment of manuscripts. It contributes one fact that bears on this
submission: its Table 1 lists a badly written or not understandable paper among
the acceptable reasons for a referee to decline an invitation, so clarity
affects whether a manuscript is reviewed at all.

\subsection{Measurable compliance}

\begin{center}
\begin{tabular}{llll}
\toprule
Limit & Journal & Manuscript & \\
\midrule
Abstract & 250 words & 238 & within \\
Body text & 8000 words & 6580 & within \\
Figures & 8 & 8 & at the limit \\
Tables & 3 & 4 & \textbf{over} \\
Keywords & 1 to 7 & 6 & within \\
Highlights & 3 to 5 bullets, 85 characters & 5, all within & within \\
\bottomrule
\end{tabular}
\end{center}

One table must go, and the smallest is not the cheapest to lose. Table 1 holds
the published source for every canopy parameter, which is what supports the claim
that nothing was fitted. Table 2, the ensemble members, is reproducible as a
sentence plus the member code already defined, and is the least destructive
reduction. Figures stand exactly at the limit, so any figure added in revision
displaces another.

\subsection{Criterion by criterion}

\begin{center}
\small
\begin{tabularx}{\textwidth}{lXl}
\toprule
Criterion & State & Exposure \\
\midrule
Title describes the content & Names the three factors and the ranking & none \\
Abstract stands alone, with limitations & Reports the ranking with intervals; carries no limitation & moderate \\
Introduction current, objectives at the end & 22 of 29 references predate 2015; closes on a roadmap paragraph & high (item 3) \\
Methods replicable & Nine settings ungrounded; no convergence test & high (item 5) \\
Results well presented & Four tables against three; a deployment-level sensitivity exists and is not reported & moderate \\
Conclusions supported, limitations emphasised & A Limitations subsection names the uncalibrated erosion threshold & low \\
References up to date & 22 of 29 predate 2015 & high (item 3) \\
Language quality & No professional editing & moderate \\
\bottomrule
\end{tabularx}
\end{center}

\subsection{Scope}

The Guide for Authors states that the journal does not accept manuscripts on
data of mainly local interest, and discourages reports of mainly local and
descriptive nature. This is the one criterion on which a rejection would be
structural rather than correctable in revision, and it is applied by the
editors before referees are invited, on the title, abstract and Introduction
alone.

The manuscript's defence is that its claims concern what a class of validation
data can certify, and how a factorial behaves when a dominant factor is
omitted, neither of which is specific to this basin. That defence currently
sits in Sections 5 and 6. Removing the geometry from the Introduction, as
observation 2 proposes, acts in the same direction.

\subsection{Points arising from the criteria}

\begin{enumerate}[leftmargin=1.5em, itemsep=0.15em]
  \item Four tables against a limit of three.
  \item The abstract carries no limitation.
  \item The Introduction closes on a roadmap rather than on the objectives.
  \item The generality claim is absent from the abstract and the Introduction.
  \item Language editing before submission.
  \item A deployment-level sensitivity is computed and not reported.
\end{enumerate}

\vspace{0.4em}
{\footnotesize\itshape Dipartimento di Ingegneria, Universit\`a degli Studi di
Palermo \textperiodcentered\
\href{mailto:cicero.martinsjr@unipa.it}{cicero.martinsjr@unipa.it}}

\end{document}
